\documentclass[11pt]{article}

\usepackage[margin=1in]{geometry}
\usepackage{lmodern}
\usepackage[T1]{fontenc}
\usepackage{microtype}
\usepackage{amsmath, amssymb, amsthm, mathtools, bm}
\usepackage{booktabs}
\usepackage{enumitem}
\usepackage{xcolor}
\usepackage{float}
\usepackage{algorithm}
\usepackage{algpseudocode}
\usepackage[colorlinks=true, linkcolor=blue!50!black, citecolor=blue!50!black,
            urlcolor=blue!50!black]{hyperref}

\newtheorem{proposition}{Proposition}
\newtheorem{lemma}{Lemma}
\theoremstyle{remark}
\newtheorem{remark}{Remark}

% ---- notation -----------------------------------------------------------------------------
\newcommand{\x}{\bm{x}}
\newcommand{\y}{\bm{y}}
\newcommand{\uu}{\bm{u}}
\newcommand{\eps}{\bm{\varepsilon}}
\newcommand{\zz}{\bm{z}}
\newcommand{\E}{\mathbb{E}}
\newcommand{\C}{\mathbb{C}}
\newcommand{\CN}{\mathcal{CN}}
\newcommand{\Id}{\bm{I}}
\newcommand{\Op}{\bm{E}}
\newcommand{\D}{D_\theta}
\newcommand{\ssb}{\sigma_{\mathrm{sb}}}
\newcommand{\xh}{\hat{\x}_0}
\newcommand{\epsh}{\hat{\eps}}
\newcommand{\grad}{\nabla}
\newcommand{\code}[1]{\texttt{#1}}

\title{I$^2$SB through the ImMAP lens:\\[2pt]
       \large predict the endpoint, then resample on the bridge}
\author{ImMAP project notes}
\date{November 2026}

\begin{document}
\maketitle

\begin{abstract}
We give a two-line reading of the I$^2$SB sampler. The trained regressor computes the
conditional expectation of the bridge's endpoint, $\xh = \D(\x_t,t,\x_1) =
\E[\x_0\mid\x_t,\x_1]$. One reverse step of I$^2$SB is then exactly: \emph{(i)} predict the
endpoint, and \emph{(ii)} resample on the bridge between that prediction and $\x_1$ at a time
$s<t$ closer to the target. The resample commits to some randomness; the next prediction
conditions on it and is sharper. Along the exact reverse bridge the predictions form a
martingale, from the blurry posterior mean $\E[\x_0\mid\x_1]$ at $t=1$ to a single sample of
$\x_0$ at $t=0$. The only detail I$^2$SB adds is that the resample re-uses a fraction
$r = \sigma(s)/\sigma(t)$ of the noise implied by the prediction; unrolled, the I$^2$SB state is
a bridge sample around a precision-weighted average of all endpoint predictions so far.
Resampling with entirely fresh noise ($r=0$) is the memoryless member of the same family.
\end{abstract}

\paragraph{Notation.}
Images are vectors in $\C^N$; $\eps,\zz \sim \CN(\bm 0, \Id)$. $\x_0$ is the target (the
fully-sampled image), $\x_1$ the bridge's prior endpoint (in MRI: the zero-filled image
$\Op^H\y$, the ACS low-pass image, or a guide contrast), and $t\in[0,1]$ the bridge time,
$t=0$ at $\x_0$. Gradients of real functions of complex vectors are Wirtinger gradients
$\partial/\partial\bar{\uu}$, so all formulas read as in the real case.

%=============================================================================================
\section{The bridge, and what the network learns}
%=============================================================================================

The Brownian bridge between $\x_0$ and $\x_1$ has marginals
\begin{equation}
  \x_t = (1-t)\,\x_0 + t\,\x_1 + \ssb(t)\,\eps, \qquad \ssb(t)^2 = V\,t(1-t),
  \qquad \eps \perp (\x_0,\x_1),
  \label{eq:marginal}
\end{equation}
with total variance $V = 4\tau^2$ set by the temperature $\tau$ ($\max_t \ssb = \tau$).
I$^2$SB trains a regressor on pairs $(\x_0,\x_1)$, $t\sim\mathcal U[0,1]$, $\x_t$ drawn from
\eqref{eq:marginal}, with the squared loss $\E\|\D(\x_t,t,\x_1) - \x_0\|^2$. Its minimiser is
the conditional expectation,
\begin{equation}
  \D(\x_t,t,\x_1) = \E[\x_0\mid\x_t,\x_1].
  \label{eq:mmse}
\end{equation}

\begin{remark}[General schedules]
For an I$^2$SB schedule $\beta(t)$, $\x_t = \mu_0(t)\x_0 + \mu_1(t)\x_1 + \ssb(t)\eps$ with
$\mu_0+\mu_1 = 1$. Everything below holds with $(1-t,\,t)$ replaced by $(\mu_0,\mu_1)$.
\end{remark}

%=============================================================================================
\section{Step 1: predict the endpoint}
%=============================================================================================

At every step the network returns its prediction of the bridge's endpoint,
\begin{equation}
  \boxed{\;\xh \coloneqq \D(\x_t,t,\x_1) = \E[\x_0\mid\x_t,\x_1].\;}
  \label{eq:predict}
\end{equation}
This is not a small move toward $\x_0$: it is the full conditional expectation of the endpoint
given everything the state reveals. How good it is depends on how much the state reveals,
which the bridge time measures.

\begin{lemma}[Bridge time is a noise level]\label{lem:debias}
Fix $\x_1$ and $t<1$, and set $\uu_t \coloneqq (\x_t - t\,\x_1)/(1-t)$. Then
\begin{equation}
  \uu_t = \x_0 + \sigma(t)\,\eps,\qquad
  \sigma(t)^2 = \frac{\ssb(t)^2}{(1-t)^2} = \frac{V t}{1-t},
  \label{eq:denoising}
\end{equation}
and $t\mapsto\sigma(t)$ is a bijection $[0,1)\to[0,\infty)$.
\end{lemma}
\begin{proof}
Substitute \eqref{eq:marginal}.
\end{proof}

So given $\x_1$, the state at time $t$ is the target plus white noise of std $\sigma(t)$, and
since $(\x_t,\x_1)\mapsto(\uu_t,\x_1)$ is invertible, $\xh = \E[\x_0\mid\uu_t,\x_1]$ is the
conditional MMSE denoiser at level $\sigma(t)$. At $t=1$ ($\sigma=\infty$) the state carries
no information beyond $\x_1$, and $\xh = \E[\x_0\mid\x_1]$: the posterior mean over every
plausible target, which is blurry. As $t\to0$ the noise vanishes and $\xh\to\x_0$.

\begin{remark}[Tweedie's formula]
By Tweedie's formula (Miyasawa), the same prediction can be read as a gradient step on the
noise-smoothed conditional prior $p_\sigma(\uu\mid\x_1) = \int\varphi_\sigma(\uu-\x)\,
p(\x\mid\x_1)\,d\x$:
\[
  \xh = \uu_t + \sigma(t)^2\,\grad_{\uu}\log p_{\sigma(t)}(\uu_t\mid\x_1).
\]
This is an identity, not an extra computation: the network computes the expectation directly,
and the formula reinterprets it as a full jump of step size $\sigma^2$ up the smoothed
log-density. The reading below does not need it.
\end{remark}

%=============================================================================================
\section{Step 2: resample on the bridge, closer to the target}
%=============================================================================================

We do not stop at $\xh$: at large $\sigma(t)$ it is an average over many targets, not one of
them. Instead we treat it as a provisional endpoint and draw a new state on the bridge
between $\xh$ and $\x_1$, at a time $s<t$, i.e.\ at a smaller noise level
$\sigma(s)<\sigma(t)$:
\begin{equation}
  \x_s = (1-s)\,\xh + s\,\x_1 + \ssb(s)\,\eps_s .
  \label{eq:resample}
\end{equation}
The noise $\eps_s$ can be drawn fresh, or partly re-used from the current state. The noise
that the prediction implies is in $\x_t$ is
\begin{equation}
  \epsh_t \coloneqq \frac{\x_t - (1-t)\,\xh - t\,\x_1}{\ssb(t)} = \frac{\uu_t - \xh}{\sigma(t)},
  \label{eq:implied_noise}
\end{equation}
and we take
\begin{equation}
  \eps_s = r\,\epsh_t + \sqrt{1-r^2}\;\zz, \qquad \zz\sim\CN(\bm0,\Id)\ \text{fresh},
  \qquad r\in[0,1].
  \label{eq:noise_mix}
\end{equation}
We write $\x_s \sim \mathcal R_s(\xh;\,\x_t)$ for this resample, \eqref{eq:resample} with
\eqref{eq:noise_mix}. If $\xh = \x_0$ then $\epsh_t = \eps$ is standard normal, $\eps_s$ is too, and
\eqref{eq:resample} is an exact draw from the bridge marginal at $s$, the input distribution
the network was trained on.

The resample is where the sampler commits: the fresh noise $\zz$ picks out, at the coarse
scale $\sigma(s)$, one of the targets that $\xh$ was averaging over. The next prediction
conditions on that choice.

%=============================================================================================
\section{The predictions sharpen: a martingale}
%=============================================================================================

\begin{proposition}[Endpoint predictions form a martingale]\label{prop:martingale}
Let $1 = t_0 > t_1 > \dots > t_K = 0$, and let $\x_{t_0},\x_{t_1},\dots$ be a path of the
bridge \eqref{eq:marginal} visited from $\x_1$ toward $\x_0$. Then
\begin{equation}
  \xh^{(k)} \coloneqq \E[\x_0\mid\x_{t_k},\x_1]
  = \E[\x_0\mid\x_{t_k},\x_{t_{k-1}},\dots,\x_{t_0}],
  \label{eq:martingale}
\end{equation}
so $(\xh^{(k)})_k$ is a Doob martingale: $\E[\xh^{(k+1)}\mid\x_{t_k},\dots,\x_{t_0}] = \xh^{(k)}$.
It starts at the posterior mean $\xh^{(0)} = \E[\x_0\mid\x_1]$ and ends at
$\xh^{(K)} = \x_0$.
\end{proposition}
\begin{proof}
Given its endpoints, the Brownian bridge is Markov. So given $\x_{t_k}$, the states at later
times $t_{k-1},\dots,t_0$ (closer to $\x_1$) carry no further information about $\x_0$, which
gives the equality in \eqref{eq:martingale}. The right-hand side conditions on a growing
history, so the tower property gives the martingale property. At $t_K=0$ the state is $\x_0$.
\end{proof}

This is the precise sense in which each step ``predicts the endpoint, and the predictions
converge to it''. There is no fixed $\x_0$ waiting to be found at sampling time: the sampler
creates the endpoint as it goes. The first prediction is the posterior mean over all plausible
targets; each resample reveals more of one particular target; each later prediction conditions
on what has been revealed and is sharper; at $t=0$ the predictions have collapsed onto a single
sample of $p(\x_0\mid\x_1)$.

\begin{remark}[Exact only for the true reverse bridge]
Proposition~\ref{prop:martingale} holds when the states are drawn from the true reverse
kernel $p(\x_s\mid\x_t,\x_1) = \int q(\x_s\mid\x_t,\x_0)\,p(\x_0\mid\x_t,\x_1)\,d\x_0$, a
mixture over plausible endpoints. I$^2$SB substitutes the point prediction $\xh$ for $\x_0$ in
$q(\x_s\mid\x_t,\x_0)$ (Section~\ref{sec:i2sb}). That gets the mean of $\x_s$ right but
drops the spread of $\x_0$ given $\x_t$ from its variance. The gap vanishes as the steps get
small.
\end{remark}

%=============================================================================================
\section{I\texorpdfstring{$^2$}{2}SB is this, with \texorpdfstring{$r=\sigma(s)/\sigma(t)$}{r = sigma(s)/sigma(t)}}\label{sec:i2sb}
%=============================================================================================

I$^2$SB's reverse step is the bridge posterior $q(\x_s\mid\x_t,\x_0)$ with $\x_0$ replaced
by $\xh$:
\begin{equation}
  \x_s \sim \CN\!\Big(\big(1-\tfrac{s}{t}\big)\xh + \tfrac{s}{t}\,\x_t,\;
  V\,\tfrac{s(t-s)}{t}\,\Id\Big)
  \label{eq:posterior}
\end{equation}
(\code{sb.base.reverse\_sample}, \code{posterior="ddpm"}).

\begin{proposition}[I$^2$SB = endpoint prediction + bridge resample]\label{prop:main}
The resample \eqref{eq:resample}--\eqref{eq:noise_mix} with
\begin{equation}
  r = \frac{\sigma(s)}{\sigma(t)} = \sqrt{\frac{s(1-t)}{t(1-s)}}
  \label{eq:r}
\end{equation}
is exactly the I$^2$SB step \eqref{eq:posterior}.
\end{proposition}
\begin{proof}
Substitute \eqref{eq:implied_noise} into \eqref{eq:resample}. The coefficient of $\x_t$ is
$k = r\,\ssb(s)/\ssb(t)$. With $\ssb = (1-t)\sigma$ and $\sigma^2 = Vt/(1-t)$,
\[
  k = \frac{(1-s)\,\sigma(s)^2}{(1-t)\,\sigma(t)^2} = \frac{s}{t}.
\]
The coefficient of $\x_1$ is $s - k\,t = 0$, and that of $\xh$ is
$(1-s) - k(1-t) = (t-s)/t$: the mean of \eqref{eq:posterior}. The fresh-noise variance is
$\ssb(s)^2(1-r^2) = V s(1-s)\cdot\frac{t-s}{t(1-s)} = V\frac{s(t-s)}{t}$: the variance of
\eqref{eq:posterior}.
\end{proof}

The identity holds for general schedules with $\sigma = \ssb/\mu_0$. We checked it
numerically against \code{reverse\_sample} for the Brownian and I$^2$SB-$\beta$ schedules
(agreement to $10^{-14}$).

\paragraph{When the re-used noise matters.}
$r$ depends only on how far the step moves.
\begin{itemize}[leftmargin=1.5em]
  \item \textbf{First step from $\x_1$:} $\sigma(1) = \infty$, so $r=0$. I$^2$SB's first step is
        an endpoint prediction followed by a fully fresh resample.
  \item \textbf{Last step to $s=0$:} $\ssb(0)=0$, so the output is $\xh$ whatever $r$ is.
  \item \textbf{Coarse ladders} (few NFE): $r$ is small (e.g.\ $0.27$ for $t=0.6\to s=0.1$), and
        I$^2$SB is nearly a fresh resample.
  \item \textbf{Fine ladders:} $r\to1$; most of the noise is carried over and the step moves the
        state only slightly toward $\xh$.
\end{itemize}

%=============================================================================================
\section{What the re-used noise does: averaging the predictions}
%=============================================================================================

In debiased coordinates the step is especially simple. With $\uu_s = \xh + \sigma(s)\eps_s$
and \eqref{eq:implied_noise}--\eqref{eq:noise_mix},
\begin{equation}
  \uu_s = r^2\,\uu_t + (1-r^2)\,\xh + \sigma(s)\sqrt{1-r^2}\;\zz .
  \label{eq:u_step}
\end{equation}
So I$^2$SB moves the state only a fraction $1-r^2 = 1 - \sigma(s)^2/\sigma(t)^2$ of the way
to the current prediction, and re-noises to $\sigma(s)$.

\begin{proposition}[Precision-weighted average]\label{prop:average}
Run \eqref{eq:u_step} down a ladder $\sigma_0=\infty > \sigma_1 > \dots > \sigma_K$, with
endpoint predictions $\xh^{(k)}$ made at level $\sigma_k$. Then
\begin{equation}
  \uu_K = \bar{\uu}_K + \bm n_K,\qquad
  \bar{\uu}_K = \sigma_K^2\sum_{k=0}^{K-1}\Big(\frac{1}{\sigma_{k+1}^2} - \frac{1}{\sigma_k^2}\Big)\,\xh^{(k)},
  \label{eq:average}
\end{equation}
where the weights are nonnegative and sum to one, and $\bm n_K$ is the sampler's own noise,
of std $\sigma_K$.
\end{proposition}
\begin{proof}
Split $\uu_k = \bar\uu_k + \bm n_k$ into the part driven by the $\xh$'s and the part driven
by the injected noise. Dividing the $\bar\uu$ part of \eqref{eq:u_step} by $\sigma_{k+1}^2$
and using $r^2 = \sigma_{k+1}^2/\sigma_k^2$ gives
$\bar\uu_{k+1}/\sigma_{k+1}^2 = \bar\uu_k/\sigma_k^2 + (\sigma_{k+1}^{-2} - \sigma_k^{-2})\,\xh^{(k)}$,
which telescopes from $1/\sigma_0^2 = 0$. The weights sum to
$\sigma_K^2(\sigma_K^{-2} - \sigma_0^{-2}) = 1$. For the noise,
$\operatorname{Var}\bm n_{k+1} = r^4\sigma_k^2 + \sigma_{k+1}^2(1-r^2) = \sigma_{k+1}^2$.
\end{proof}

This is the whole picture of I$^2$SB: \emph{at each step, predict the endpoint; then resample
on the bridge, at a lower noise level, around a running average of all endpoint predictions so
far, each weighted by the precision $1/\sigma^2$ gained when it was made.} Early, blurry
predictions get little weight; late, sharp ones get most of it. The decomposition is an
algebraic identity of the sampler; the network enters only through the $\xh^{(k)}$, each of
which depends on the state it was computed from.

\begin{remark}[Fresh resampling, $r=0$]
Drawing $\eps_s$ entirely fresh gives $\x_s = (1-s)\xh + s\x_1 + \ssb(s)\zz$, the forward bridge
sample around the latest prediction (\code{posterior="interpolant"}). It is memoryless: the
state is centred on the last prediction only, and nothing is averaged. With a perfect network,
fresh resampling and I$^2$SB both produce exact bridge marginals; they differ only in how the
network's errors propagate, which makes \code{interpolant} vs.\ \code{ddpm} a direct test of
the value of the averaging.
\end{remark}

\begin{algorithm}[H]
\caption{I$^2$SB as endpoint prediction + bridge resample}
\label{alg:i2sb}
\begin{algorithmic}[1]
\Require $\x_1$; regressor $\D$; times $1 = t_0 > t_1 > \dots > t_K = 0$
\State $\x \gets \x_1$
\For{$k = 0,\dots,K-1$}
  \State $\xh \gets \D(\x, t_k, \x_1)$
         \Comment{predict the endpoint}
  \State $\x \sim \mathcal R_{t_{k+1}}(\xh;\,\x)$
         \Comment{resample on the bridge}
\EndFor
\State \Return $\x$
\end{algorithmic}
\end{algorithm}
\noindent Here $\mathcal R$ uses $r_k = \sigma(t_{k+1})/\sigma(t_k)$ for I$^2$SB and $r_k = 0$
for fresh resampling. The last step lands at $t_K = 0$, where $\ssb = 0$, so the returned
state is the final prediction $\xh$.

%=============================================================================================
\section{Summary}
%=============================================================================================

\begin{table}[h]
\centering
\small
\begin{tabular}{@{}lll@{}}
\toprule
Object & Formula & Code \\
\midrule
bridge marginal & $\x_t = (1-t)\x_0 + t\x_1 + \ssb(t)\eps$, $\ssb^2 = Vt(1-t)$ &
  \code{sb.base.bridge\_coeffs} \\
time $\leftrightarrow$ noise & $\uu_t = \x_0 + \sigma(t)\eps$, $\sigma^2 = Vt/(1-t)$ & --- \\
endpoint prediction & $\xh = \D(\x_t,t,\x_1) = \E[\x_0\mid\x_t,\x_1]$ & \code{predict\_x0} \\
bridge resample & $\x_s = (1-s)\xh + s\x_1 + \ssb(s)\eps_s$ & \code{reverse\_sample} \\
noise re-use & $\eps_s = r\epsh_t + \sqrt{1-r^2}\zz$ & \\
\quad I$^2$SB & $r = \sigma(s)/\sigma(t)$ & \code{posterior="ddpm"} \\
\quad fresh & $r = 0$ & \code{posterior="interpolant"} \\
predictions & martingale, $\E[\x_0\mid\x_1]\to\x_0$ (exact reverse bridge) & --- \\
I$^2$SB state & precision-weighted average of the $\xh^{(k)}$, plus noise & --- \\
\bottomrule
\end{tabular}
\caption{I$^2$SB as endpoint prediction followed by a bridge resample (\code{sb/base.py}).}
\label{tab:code}
\end{table}

\begin{thebibliography}{9}
\bibitem{liu2023}
G.-H.~Liu, A.~Vahdat, D.-A.~Huang, E.~A.~Theodorou, W.~Nie, A.~Anandkumar.
I$^2$SB: Image-to-Image Schr\"odinger Bridge. \emph{ICML}, 2023.

\bibitem{efron2011}
B.~Efron. Tweedie's formula and selection bias.
\emph{Journal of the American Statistical Association}, 106(496), 2011.

\bibitem{miyasawa1961}
K.~Miyasawa. An empirical Bayes estimator of the mean of a normal population.
\emph{Bulletin of the International Statistical Institute}, 38, 1961.
\end{thebibliography}

\end{document}
